% !TEX program = pdflatex
% Erdős #132, cases 7 <= n <= 13 — integrated draft (2026-07-03, evening pass)
\documentclass[11pt]{article}

\usepackage[margin=1.1in]{geometry}
\usepackage{amsmath,amssymb,amsthm}
\usepackage{booktabs}
\usepackage{xcolor}
\usepackage[colorlinks=true,linkcolor=blue!60!black,citecolor=blue!60!black,urlcolor=blue!60!black]{hyperref}

% ---- todo marker (self-contained; no todonotes dependency) ----
\newcommand{\todo}[1]{\textcolor{red!70!black}{\textbf{[TODO: #1]}}}

\newtheorem{theorem}{Theorem}[section]
\newtheorem{lemma}[theorem]{Lemma}
\newtheorem{corollary}[theorem]{Corollary}
\newtheorem{proposition}[theorem]{Proposition}
\theoremstyle{definition}
\newtheorem{definition}[theorem]{Definition}
\theoremstyle{remark}
\newtheorem{remark}[theorem]{Remark}

\newcommand{\R}{\mathbb{R}}
\newcommand{\Q}{\mathbb{Q}}
\newcommand{\Z}{\mathbb{Z}}
\newcommand{\Lat}{L_{\Delta}}
\DeclareMathOperator{\diam}{diam}

\title{On distances of low multiplicity:\\ the cases $7 \le n \le 13$ of a problem of Erd\H{o}s}
\author{Juan Patricio Marchetto\thanks{Independent researcher.}}
\date{July 2026}

\begin{document}
\maketitle

\begin{abstract}
Let $A \subset \R^2$ be a set of $n \ge 5$ points. Erd\H{o}s asked whether there
must be at least two distances that occur at least once and between at most $n$ pairs of
points (Erd\H{o}s problem \#132); Hopf and Pannwitz proved that the diameter is always such
a distance. Erd\H{o}s and Fishburn settled $n = 5, 6$; for unrestricted point sets every
case $n \ge 7$ had remained open. We combine an elementary counting reduction (which
forces at most $\lfloor n/2 \rfloor$ distinct distances in any counterexample, with a
rigid multiplicity profile when $n$ is even) with a descent lemma along the diameter
graph and with the published classifications of planar few-distance sets
(Erd\H{o}s--Fishburn; Shinohara; Wei; Sz\"oll\H{o}si--\"Osterg\aa rd) to prove the
conjecture for $n \in \{7, 8, 9, 10, 13\}$ unconditionally, and for $n = 11$ and $n = 12$
conditionally on Wei's classification of $11$-point $5$-distance sets \cite{Wei11}: the
sole nontrivial branch of $n = 12$ (a rigid six-distance profile) reduces by descent to a
finite exact extension search over a single explicitly described $11$-point configuration
(the regular $11$-gon), which we have carried out in exact arithmetic in the field
$\Q(\zeta_{11})^+$, with no solution. The case $n = 10$ is closed by the same method ---
descent to the classified $9$-point $4$-distance sets of Erd\H{o}s--Fishburn, then an exact
extension search over the regular nonagon in $\Q(\zeta_9)^+$ with no solution --- so it
depends on no classification beyond the fully-proved $E_9(4)$. The entire range
$7 \le n \le 13$ is thus settled;
the only input that is not self-contained is the classification $E_{11}(5)$ of
\cite{Wei11}, on which the cases $n = 11$ and $n = 12$ depend and whose sole published
source omits several of its proofs for length; we therefore accompany them with an
independent exhaustive re-derivation, in exact arithmetic, of the triangular-lattice part
of that classification. The counting reduction,
the case $n = 7$, and the reduction of $n = 8$ to the profile $(1,9,9,9)$ are due to
Zeraoulia; the relevance of the $8$-point $4$-distance classification to $n = 8$ was first
observed in a forum comment by Chojecki. As by-products we record an erratum to Wei's
classification of $10$-point $5$-distance sets (two of the twenty listed configurations
are similar, so the list has $19$ classes), to our knowledge the first published
multiplicity tables for the complete $8$-point $4$-distance classification, and exact
multiplicity profiles for
the four classes of $E_{11}(5)$ together with the independent re-derivation of its three
triangular-lattice classes. All finite verifications are carried out in exact arithmetic
and are reproducible from the ancillary files.
\end{abstract}

\section{Introduction}\label{sec:intro}

For a finite set $A \subset \R^2$ of $n$ points, call a distance $\delta$ realized by $A$
\emph{$n$-good} if $1 \le m(\delta) \le n$, where $m(\delta)$ denotes the number of pairs
$\{x,y\} \subseteq A$ with $|xy| = \delta$. Erd\H{o}s \cite{Er84c} asked whether a
second such distance must always exist; Erd\H{o}s and Pach \cite{ErPa90} asked further
whether the number of such distances must tend to infinity
(see also \cite{ErFi95,Bloom132}):

\begin{quote}
\emph{Let $A \subset \R^2$ be a set of $n$ points. Must there be two distances which occur
at least once but between at most $n$ pairs of points? Must the number of such distances
tend to infinity as $n \to \infty$?}
\end{quote}

Hopf and Pannwitz \cite{HoPa34} proved that the largest distance (the diameter) occurs at
most $n$ times, so the diameter is always $n$-good; the question is whether a
\emph{second} $n$-good distance must exist. Erd\H{o}s \cite{Er84c} believed this holds for
all $n \ge 5$; it fails for $n = 4$: two unit equilateral triangles glued along an edge (a
rhombus with angles $60^\circ$ and $120^\circ$) realize the distance $1$ on five pairs
and the long diagonal $\sqrt{3}$ on one pair, so the multiplicity vector, ordered by
decreasing distance, is $(1,5)$, and the diameter is the \emph{only} $4$-good distance
($5 > 4$). Erd\H{o}s and Fishburn \cite{ErFi95} proved the statement for $n = 5$ and
$n = 6$. Recently Clemen, Dumitrescu and Liu \cite{CDL25} proved it for all point sets in
convex position (and for ``not too convex'' sets); for unrestricted point sets every case
$n \ge 7$ has remained open. Erd\H{o}s offered \$100 for ``any nontrivial result''
\cite{Er97e}. For the state of the problem see the curated page of Bloom \cite{Bloom132}.

Our main result is the following.

\begin{theorem}[main]\label{thm:main}
Let $A \subset \R^2$ be a set of $n$ points.
\begin{enumerate}
\item[(i)] If $n \in \{7, 8, 9, 10, 13\}$, then $A$ has at least two $n$-good distances.
\item[(ii)] If $n = 12$, then $A$ has at least two $12$-good distances, conditionally on
the classification $E_{11}(5)$ of $11$-point $5$-distance sets \cite[Theorem~13]{Wei11}: a
failure would force exactly six distinct distances with the rigid multiplicity vector,
ordered by decreasing distance, $(1, 13, 13, 13, 13, 13)$; by descent this branch is
equivalent to adjoining one point to a single $11$-point configuration, the regular
$11$-gon $R_{11}$ (the only class in the relevant window, forcing the degree pattern
$a = (2,2,2,2,2)$), and the resulting finite exact extension search over $R_{11}$ in
$\Q(\zeta_{11})^+$ has no solution (zero survivors); hence no such $A$ exists (Corollary
\ref{cor:extension}, Proposition \ref{prop:n12search} and Section \ref{sec:n12}).
\item[(iii)] If $n = 11$, then $A$ has at least two $11$-good distances, conditionally on
the classification $E_{11}(5)$ of $11$-point $5$-distance sets \cite[Theorem~13]{Wei11}:
a failure would reduce to five distinct distances with diameter multiplicity in
$\{2, \dots, 7\}$ and a non-star diameter graph, hence to an $11$-point $5$-distance set,
and none of the four similarity classes of $E_{11}(5)$ has such a multiplicity profile
(Section \ref{sec:n11}).
\end{enumerate}
\end{theorem}

Both parts (ii) and (iii) rest on the classification $E_{11}(5)$ of $11$-point
$5$-distance sets due to Wei \cite{Wei11} (Theorem~13: four similarity classes ---
the regular $11$-gon $R_{11}$ and three triangular-lattice patches). We have obtained the
primary source (a scan of the printed paper); it establishes $E_{11}(5)$ but omits several
of its proofs for length (Section \ref{sec:n11}, Remark \ref{rem:weiomit}), so we treat
$n = 11$ and $n = 12$ as closed \emph{conditionally} on \cite{Wei11} and independently
re-derive the triangular-lattice part of the classification in exact arithmetic. Given
$E_{11}(5)$, part (iii) is a finite profile check (Section \ref{sec:n11}), and part (ii)
is a finite exact extension search (Corollary \ref{cor:extension} and Section
\ref{sec:n12}) which the same profile analysis reduces to the single candidate $R_{11}$;
we have executed that search in exact arithmetic over $\Q(\zeta_{11})^+$ (the descent
computation \texttt{e12\_descent} of \texttt{rust/fewdist}, Section \ref{sec:n12}), and it
returns no extension --- zero survivors --- so the branch is empty and $n = 12$ is closed.
Consequently the entire range $7 \le n \le 13$ is now settled; nothing in it remains open
by computation. The one ingredient that is not self-contained is the citation of
\cite{Wei11} for $E_{11}(5)$ (used by both $n = 11$ and $n = 12$), which we re-derive in
part --- its triangular-lattice classes --- but not in full.

\subsection*{Attribution and history of the cases $n = 7, 8$}

Scrupulous attribution is essential here, since partial claims on this problem were posted
publicly in January 2026.

\begin{itemize}
\item \textbf{Zeraoulia (2026).} The counting (pigeonhole) reduction, the complete proof
of the case $n = 7$, and the reduction of $n = 8$ to the rigid multiplicity profile
$(1,9,9,9)$ are due to Rafik Zeraoulia, in a self-published working draft on ResearchGate
\cite{Zeraoulia26} (January 28, 2026) and in a comment on the problem's discussion forum.
We have verified his proof of $n = 7$ step by step; it is correct and complete, resting
only on refereed literature (he quotes the classification of $7$-point $3$-distance sets
from \cite[\S 1]{Wei12}; the primary source is \cite[Theorem~1]{ErFi96}, see Remark
\ref{rem:E73}). His note is not on arXiv and has not been refereed; we reproduce the
statements we use with proofs, and none of Sections
\ref{sec:counting}--\ref{sec:n8} up to and including the profile reduction for $n=8$
should be credited to us.
\item \textbf{Chojecki (2026).} Three hours after Zeraoulia's comment, Przemek Chojecki
posted on the same forum thread \cite{Chojecki26} an output of the language model GPT-5.2
sketching a closure of $n = 8$ via ``the classification of planar 8-point 4-distance
sets''. That classification does exist --- it is Theorem 1.2(a) of Shinohara's
uniqueness paper \cite{Shinohara08} --- but the posted sketch cites no source, exhibits no
multiplicity tables, and has no author of record; we cite it as a communication for the
\emph{observation} that \cite{Shinohara08} is the relevant reference, and we supply the
verification independently. Remark \ref{rem:gpt} records the exact comparison between
the posted sketch and the paper.
\item \textbf{Classification literature.} The finite-case analysis below consumes the
classifications of planar $k$-distance sets: the $7$-point $3$-distance sets and the
$9$-point $4$-distance sets, together with the values $g(k)$ for $k \le 5$
(Erd\H{o}s--Fishburn \cite{ErFi96}); all planar $3$-distance sets (Shinohara
\cite{Shinohara04}); the $7$-point $4$-distance sets (Lan--Wei \cite{LanWei13}); the
$8$-point $4$-distance sets and the uniqueness of the maximum $12$-point $5$-distance
set (Shinohara \cite{Shinohara08}); the value $g(6) = 13$, together with the $10$-point
$5$-distance sets used for $n = 10$ only as an independent corroboration (Remark
\ref{rem:n10wei}), the main proof for $n = 10$ descending instead to $E_9(4)$
(Wei \cite{Wei12}); the $11$-point $5$-distance sets (Wei \cite{Wei11}); and the
$13$-point $6$-distance sets (Sz\"oll\H{o}si--\"Osterg\aa rd \cite{SzOs20}, agreeing
with \cite{Wei12}). Every classification statement used below has been checked against the
full text of its primary source, and every multiplicity profile has been recomputed in
exact arithmetic; the one caveat is \cite{Wei11}, whose printed text (a scan, now in hand)
states the classification $E_{11}(5)$ but omits several of its proofs for length (``some
proofs are omitted because of the restriction of the length of the paper'',
\cite[p.~506]{Wei11}). We therefore re-derive its triangular-lattice part independently
and flag the case $n = 11$ as conditional on \cite{Wei11} (Section \ref{sec:n11}).
\end{itemize}

\subsection*{Reproducibility (ancillary files)}

No decision in any finite verification below relies on floating-point arithmetic. The
computations are packaged as ancillary files (in the arXiv sense) in the directory
\texttt{tools/} accompanying this note:
\texttt{n8\_descent.py} (the exact descent search that closes $n = 8$, Section
\ref{sec:n8}; pure standard library, with an optional independent re-verification layer
on top of \texttt{sympy});
\texttt{n9\_profiles.py} (exact multiplicity profiles of the four $9$-point $4$-distance
sets, Section \ref{sec:n9});
\texttt{n13\_profiles.py} (exact profiles of the three $13$-point $6$-distance sets,
Section \ref{sec:n13});
\texttt{extract\_figs.py}, \texttt{match\_figs.py} and \texttt{nonlattice\_exact.py}
(extraction of the configurations of \cite[Figure~2]{Wei12}, matching against an
independent exhaustive re-derivation of the triangular-lattice classes, and exact
profiles of the non-lattice classes over $\Q(\sqrt5)$, Section \ref{sec:n10});
\texttt{n11\_profiles.py} (exact multiplicity profiles of the four $11$-point $5$-distance
sets of \cite[Theorem~13]{Wei11} and the window checks that close $n = 11$ and locate the
$n = 12$ candidate, Sections \ref{sec:n11}--\ref{sec:n12});
together with the reference outputs \texttt{figs.json}, \texttt{e10\_M100.txt},
\texttt{e10\_brute\_output.txt}, \texttt{e11\_lattice.json} (the exhaustive
triangular-lattice re-derivation of the three lattice classes of $E_{11}(5)$, produced by
the standalone crate \texttt{wei-xcheck} described below) and
\texttt{eleven\_point\_five\_distance.json} (the exact coordinates and multiplicity
profiles of the four classes of $E_{11}(5)$ used in Section \ref{sec:n11}).
Finally, the exact extension searches are packaged as the Rust crate \texttt{rust/fewdist}
(module \texttt{extend12.rs}): the bin \texttt{e12\_descent -{}-n10} runs the unconditional
$n = 10$ search over the regular nonagon $R_9$ in $\Q(\zeta_9)^+ = \Q(\beta)$,
$\beta = 2\cos(2\pi/9)$ (minimal polynomial $t^3 - 3t + 1$; Section \ref{sec:n10},
Proposition \ref{prop:n10search}; $855$ closed-Cayley--Menger prunings, $9$ $z$-candidates,
zero survivors), and the bin \texttt{e12\_descent -{}-r11} runs the search that closes the
six-distance branch of $n = 12$ over $R_{11}$ in $\Q(\alpha)$, $\alpha = 2\cos(2\pi/11)$
(Section \ref{sec:n12}; $2310$ prunings, zero survivors); all arithmetic is exact in the
respective field and its quadratic tower. Separately, the standalone crate
\texttt{wei-xcheck} re-derives
the triangular-lattice classes of $E_{11}(5)$ exhaustively in exact Loeschian arithmetic
(output \texttt{e11\_lattice.json}).

\subsection*{Notation}

$P \subset \R^2$, $|P| = n \ge 5$, points distinct. The distinct distances of $P$ are
$\delta_1 > \delta_2 > \dots > \delta_k$ ($\delta_1 = \diam P$), with multiplicities
$m_i = \#\{\{x,y\} \subseteq P : |xy| = \delta_i\}$ and $d := m_1$. Thus
$\sum_{i=1}^k m_i = \binom{n}{2}$. Multiplicity vectors are always ordered by
\emph{decreasing} distance, so the first entry is the diameter multiplicity. Let $g(k)$
denote the maximum size of a planar set with exactly $k$ distinct distances; we use
$g(1)=3$, $g(2)=5$, $g(3)=7$, $g(4)=9$, $g(5)=12$ \cite{ErFi96} and $g(6)=13$
\cite{Wei12}. All of $g(2), \dots, g(5)$ are \emph{proved} in \cite{ErFi96} (Theorem 1;
$g(4)=9$ in their Section 3, $g(5)=12$ in their Section 4 --- we have verified this
against the full text; what that paper leaves open is only the uniqueness of the
$12$-point $5$-distance set, later proved in \cite{Shinohara08}, and $g(6)$, later
proved in \cite{Wei12}). $\Lat$ denotes the triangular lattice, embedded by
$(a,b) \mapsto a(1,0) + b(\tfrac12, \tfrac{\sqrt3}{2})$, with squared distances the
Loeschian numbers $a^2+ab+b^2$.

\begin{definition}[failure profile]\label{def:failure}
$P$ \emph{fails} the conjecture for $n$ if it has at most one $n$-good distance. Since the
diameter is always $n$-good \cite{HoPa34}, this holds if and only if
\[
m_i \ge n+1 \quad \text{for all } i \in \{2, \dots, k\}, \qquad 1 \le d \le n .
\]
(For $n \ge 5$ we have $k \ge 2$ because $g(1) = 3$, so the equivalence is exact.)
\end{definition}

\section{The counting reduction}\label{sec:counting}

The following lemma, in the special cases $n = 7, 8$, is due to Zeraoulia
\cite{Zeraoulia26}; the general form and the parity rigidity (iii) appear to be new.

\begin{lemma}[counting reduction]\label{lem:counting}
Let $P$, $|P| = n \ge 5$, fail the conjecture. Then:
\begin{enumerate}
\item[(i)] $d + \sum_{i=2}^k m_i = \binom{n}{2}$, with $1 \le d \le n$ and $m_i \ge n+1$
for $i \ge 2$.
\item[(ii)] $(k-1)(n+1) \le \binom{n}{2} - d \le \binom{n}{2} - 1 = \frac{(n+1)(n-2)}{2}$,
and consequently $k \le \lfloor n/2 \rfloor$.
\item[(iii)] (parity rigidity) If $n$ is even and $k = n/2$, then $d = 1$ and $m_i = n+1$
for all $i \ge 2$; the multiplicity vector is exactly $(1, n+1, \dots, n+1)$.
\item[(iv)] In general $d \le \binom{n}{2} - (k-1)(n+1)$ and the vector is
$(d,\, n+1+e_2, \dots, n+1+e_k)$ with $e_i \ge 0$ and
$\sum_i e_i = \binom{n}{2} - d - (k-1)(n+1)$.
\item[(v)] $k \ge \kappa(n) := \min\{k : g(k) \ge n\}$.
\end{enumerate}
\end{lemma}

\begin{proof}
(i) Partition the $\binom{n}{2}$ pairs into distance classes; $d \le n$ is Hopf--Pannwitz
\cite{HoPa34}; $m_i \ge n+1$ is the negation of the conjecture (Definition
\ref{def:failure}); $d \ge 1$ since the diameter is realized.

(ii) From (i), $\sum_{i=2}^{k} m_i = \binom{n}{2} - d$ is a sum of $k-1$ terms each
$\ge n+1$, whence $(k-1)(n+1) \le \binom{n}{2} - d \le \binom{n}{2} - 1$. The algebraic
identity $\binom{n}{2} - 1 = \frac{n^2-n-2}{2} = \frac{(n+1)(n-2)}{2}$ gives
$k-1 \le \frac{n-2}{2}$. If $n$ is even this reads $k \le n/2$; if $n$ is odd,
$\frac{n-2}{2}$ is not an integer while $k - 1$ is, so $k - 1 \le \frac{n-3}{2}$, i.e.\
$k \le \frac{n-1}{2}$. In both cases $k \le \lfloor n/2 \rfloor$.

(iii) If $n$ is even and $k = n/2$, then
$(k-1)(n+1) = \frac{(n-2)(n+1)}{2} = \binom{n}{2} - 1$, so by (ii),
$\binom{n}{2} - d \ge \binom{n}{2} - 1$ forces $d = 1$; then
$\sum_{i \ge 2} m_i = \binom{n}{2} - 1 = (k-1)(n+1)$ with every summand $\ge n+1$ forces
term-by-term equality.

(iv) Immediate rewriting of (i) with $e_i := m_i - (n+1) \ge 0$.

(v) If $k < \kappa(n)$ then $g(k) < n = |P|$ and $P$ would be a $k$-distance set larger
than the maximum, a contradiction. Note this uses the external values of $g$, not the
counting.
\end{proof}

\begin{table}[ht]
\centering
\caption{Forced parameters of a hypothetical failure $P$, $7 \le n \le 13$. ``$k$
forced'' uses Lemma \ref{lem:counting}(ii),(v) together with the values of $g(k)$;
the ranges of $d$ additionally use Corollary \ref{cor:size} (star exclusion) where marked
$d \ge 2$. All arithmetic verified mechanically.}
\label{tab:profiles}
\begin{tabular}{@{}clll@{}}
\toprule
$n$ & $\binom{n}{2}$; fail $\Rightarrow m_i \ge$ & $k$ forced & multiplicity vector \\
\midrule
$7$ & $21$; $m_i \ge 8$ & $k = 3$ & $(d, m_2, m_3)$, $d \in \{2,\dots,5\}$ (10 vectors) \\
$8$ & $28$; $m_i \ge 9$ & $k = 4$ & rigid: $(1, 9, 9, 9)$ \\
$9$ & $36$; $m_i \ge 10$ & $k = 4$ & $(d, 10{+}e_2, 10{+}e_3, 10{+}e_4)$, $d \in \{2,\dots,6\}$ (35 vectors) \\
$10$ & $45$; $m_i \ge 11$ & $k = 5$ & rigid: $(1, 11, 11, 11, 11)$ \\
$11$ & $55$; $m_i \ge 12$ & $k = 5$ & $(d, 12{+}e_2, \dots, 12{+}e_5)$, $d \in \{2,\dots,7\}$ (126 vectors) \\
$12$ & $66$; $m_i \ge 13$ & $k \in \{5, 6\}$ & $k{=}6$: rigid $(1, 13, 13, 13, 13, 13)$; $k{=}5$: $d \in \{2,\dots,12\}$ \\
$13$ & $78$; $m_i \ge 14$ & $k = 6$ & $(d, 14{+}e_2, \dots, 14{+}e_6)$, $d \le 8$ \\
\bottomrule
\end{tabular}
\end{table}

\begin{remark}[asymptotic limitation]
The counting only gives $k \le \lfloor n/2 \rfloor$ under failure, whereas the minimum
number of distinct distances of $n$ planar points grows like $n/\log n$ up to constants
\cite{GuthKatz15}, far below $n/2$. The reduction is therefore an instrument exclusively
for small $n$; it says nothing about the general conjecture or the $\to \infty$ version.
\end{remark}

\section{Descent along the diameter graph}\label{sec:descent}

\begin{lemma}[descent]\label{lem:descent}
Let $P$, $|P| = n$, have $k \ge 2$ distinct distances with multiplicity vector
$(d, m_2, \dots, m_k)$. Assume:
\begin{itemize}
\item[\emph{(H1)}] $m_i \ge n$ for all $i \ge 2$ (under the failure profile the stronger
$m_i \ge n+1$ holds);
\item[\emph{(H2)}] some $v \in P$ is incident to \emph{all} diameter edges (equivalently,
the diameter graph is a star $K_{1,d}$ with center $v$; if $d = 1$ either endpoint
serves).
\end{itemize}
Let $P' = P \setminus \{v\}$. Then:
\begin{enumerate}
\item[(a)] $|P'| = n - 1$.
\item[(b)] The distances of $P'$ are \emph{exactly} $\{\delta_2, \dots, \delta_k\}$;
i.e.\ $P'$ is a $(k-1)$-distance set of $n-1$ points.
\item[(c)] $m_i' = m_i - \deg_i(v)$ for $i \ge 2$, where
$\deg_i(v) = \#\{p \in P' : |vp| = \delta_i\}$ and
$\sum_{i \ge 2} \deg_i(v) = n - 1 - d$. In particular
$m_i' \ge m_i - (n-1-d) \ge d + 1 \ge 2$ under (H1) (and $\ge d + 2 \ge 3$ under the
failure profile), and $\sum_{i \ge 2} m_i' = \binom{n-1}{2}$.
\item[(d)] $\diam P' = \delta_2$, and Hopf--Pannwitz applied to $P'$ gives
$m_2' \le n - 1$, hence $\deg_2(v) \ge m_2 - (n-1) \ge 2$ under the failure profile.
\end{enumerate}
\end{lemma}

\begin{proof}
(a) Trivial. (b) Every diameter edge is incident to $v$ by (H2), so all of them disappear
with $v$ and no pair of $P'$ realizes $\delta_1$; deleting a point creates no new
distances, so the distances of $P'$ are contained in $\{\delta_2, \dots, \delta_k\}$; that
each survives follows from (c), $m_i' > 0$. (c) $v$ is incident to $n-1$ pairs: $d$ of
them diameter edges (by (H2) all diameter edges pass through $v$, and if $d \ge 2$ they go
to $d$ distinct points), and the remaining $n-1-d$ distributed among classes
$2, \dots, k$. Hence $\deg_i(v) \le n-1-d$ for each $i \ge 2$, and
$m_i' = m_i - \deg_i(v) \ge n - (n-1-d) = d+1 \ge 2 > 0$ under (H1). The control sum:
$\sum_{i \ge 2} m_i' = \binom{n}{2} - d - (n-1-d) = \binom{n}{2} - (n-1) =
\binom{n-1}{2}$. (d) The maximum distance of $P'$ is $\delta_2$ by (b); Hopf--Pannwitz on
the $(n-1)$-point set $P'$ gives $m_2' \le n-1$, so
$\deg_2(v) = m_2 - m_2' \ge (n+1) - (n-1) = 2$ under the failure profile.
\end{proof}

\begin{remark}
Hypothesis (H1) is essential in (b): in general, deleting a vertex can also kill a
non-diameter class of low multiplicity, and $P'$ would have fewer than $k-1$ distances.
Under the failure profile every non-diameter class survives with at least $3$ edges.
\end{remark}

\begin{corollary}[size contradiction]\label{cor:size}
If moreover $n - 1 > g(k-1)$, then \emph{(H1)}$\,\wedge\,$\emph{(H2)} is impossible: $P'$
would be a $(k-1)$-distance set of $n-1 > g(k-1)$ points. Under the failure profile this
excludes every configuration whose diameter graph is a star --- in particular all
configurations with $d = 1$ --- whenever $n - 1 > g(k-1)$.
\end{corollary}

\begin{proof}
Immediate from Lemma \ref{lem:descent}(b) and the definition of $g$.
\end{proof}

\begin{corollary}[decidable extension]\label{cor:extension}
Suppose $n - 1 \le g(k-1)$ and the $(n-1)$-point $(k-1)$-distance sets are completely
classified: a finite list up to similarity, with algebraic coordinates. Then every
counterexample $P$ satisfying \emph{(H2)} is obtained by adjoining one point $v$ to a set
$K$ from the list, subject to:
\begin{enumerate}
\item the $n-1-d$ non-diameter distances from $v$ to $K$ take values among the $k-1$
distances of $K$ (finitely many assignments, at most $(k-1)^{\,n-1-d}$);
\item the remaining $d$ distances from $v$ are all equal to a new value
$\delta_1 > \delta_2(K)$;
\item the final profile satisfies $m_1 = d$ and $m_i = m_i(K) + \deg_i(v) \ge n+1$, with
$\deg_2(v) \ge 2$ by Lemma \ref{lem:descent}(d);
\item $v$ is determined (at most $2$ candidates per pair of constraints) as an
intersection of two circles with distinct centers in $K$ and radii among the distances of
$K$; all candidates and all remaining verifications live in $\overline{\Q}$.
\end{enumerate}
Hence the existence of such a $P$ is decidable by a finite exact-arithmetic search.
\end{corollary}

\begin{proof}
The logical direction is deletion: if $P$ is a counterexample satisfying (H2), then
$K := P \setminus \{v\}$ is an $(n-1)$-point $(k-1)$-distance set (Lemma
\ref{lem:descent}), hence appears in the classification up to similarity. Fix a
representative of $K$ with algebraic coordinates; then $P = K \cup \{v\}$ with $v$
satisfying constraints 1--4. Since $n - 1 - d \ge 2$ in all cases of interest
($n \ge 7$, $d \le n - 3$), there are at least two constraints of the form ``$v$ lies on
the circle of center $p \in K$ and radius $\delta_j(K)$'' with distinct centers; two
distinct circles meet in at most $2$ points, so the candidates $v$ are finite for each
radius assignment, and the assignments are finite. All remaining checks (the other
distances, the new value $\delta_1 > \delta_2$, equality of the $d$ far distances, the
multiplicity profile, $v \notin K$) are identities and inequalities over algebraic
numbers, hence decidable exactly. Similarity is harmless: the constraints are covariant
under similarity and the multiplicity profile is invariant.
\end{proof}

\begin{remark}[when (H2) fails]\label{rem:nostar}
Two diameter segments always intersect (share an endpoint or cross properly)
\cite{HoPa34}; and if in a family of edges every two share a vertex but no vertex lies in
all of them, the family is a triangle. Hence, in a counterexample where Corollary
\ref{cor:size} excludes the star (below: $n \in \{7, 9, 11\}$ and $n = 12$, $k = 5$), the
diameter graph contains either two properly crossing diameter edges or an equilateral
triangle of side $\delta_1$ (the latter only if $d \ge 3$). We record this as structural
ammunition; it is not used in the main line.
\end{remark}

\section{The case $n = 7$ (Zeraoulia)}\label{sec:n7}

\begin{theorem}[{Zeraoulia \cite{Zeraoulia26}}]\label{thm:n7}
Every set of $7$ points in the plane has at least two $7$-good distances.
\end{theorem}

\begin{proof}
Suppose $P$, $|P| = 7$, fails. By Lemma \ref{lem:counting}(ii) with $\binom{7}{2} = 21$
and $m_i \ge 8$: $(k-1) \cdot 8 \le 20$, so $k \le 3$. The case $k = 2$ is excluded by
$g(2) = 5 < 7$, and $k = 1$ by $g(1) = 3$; hence $k = 3$ exactly, and by Lemma
\ref{lem:counting}(iv), $d \le 21 - 2 \cdot 8 = 5$. Moreover, by Corollary \ref{cor:size}
with $g(2) = 5$: a star diameter graph (in particular $d = 1$) would yield a $2$-distance
set of $6 > 5$ points; hence $d \in \{2, \dots, 5\}$. (This last improvement is not needed
for the conclusion but sharpens the profile; it is not in \cite{Zeraoulia26}.)

Now $P$ itself is a $3$-distance set of $7 = g(3)$ points, hence maximum, and the
classification of $7$-point $3$-distance sets (\cite[Theorem~1]{ErFi96}: ``the only
7-point sets that determine three distances are $R_7$ and $R_6^+$''; independently
\cite{Shinohara04}, quoted in \cite[\S1]{Wei12}) gives exactly two configurations: the
regular heptagon $R_7$ and the regular hexagon with center $R_6^+$. Their multiplicity
vectors, ordered by decreasing distance, are
\[
R_7 : (7, 7, 7), \qquad R_6^+ : (3, 6, 12),
\]
($R_6^+$ with side $a$: diameter $2a$ with $3$ pairs, $\sqrt{3}\,a$ with $6$ pairs, $a$
with $6$ sides $+$ $6$ radii $= 12$ pairs; both vectors sum to $21$). In $R_7$ every class
has multiplicity $7 \le 7$; in $R_6^+$ the classes $3$ and $6$ are $\le 7$. Neither
matches the failure profile ($m_i \ge 8$ for $i \ge 2$), a contradiction.
\end{proof}

\begin{remark}[status of the input $E_7(3)$]\label{rem:E73}
The load-bearing list $E_7(3) = \{R_7, R_6^+\}$ is \emph{proved} in Erd\H{o}s--Fishburn
\cite{ErFi96} itself (Theorem 1, p.~116; proof in their Section 2, pp.~118--119), which
we have verified against the full text. Zeraoulia quotes the list from the survey
paragraph of \cite[\S1]{Wei12}; Shinohara's classification of all planar three-distance
sets \cite{Shinohara04} and \cite[Lemma~2.2(b)]{Shinohara08} provide independent
confirmations. The case $n = 7$ is therefore closed unconditionally.
\end{remark}

\section{The case $n = 8$}\label{sec:n8}

\begin{proposition}[reduction; {Zeraoulia \cite{Zeraoulia26}}]\label{prop:n8red}
If $P$, $|P| = 8$, fails the conjecture, then $P$ has exactly $4$ distinct distances with
multiplicity vector exactly $(1, 9, 9, 9)$; in particular the diameter is realized by a
unique pair.
\end{proposition}

\begin{proof}
$\binom{8}{2} = 28$ and $m_i \ge 9$ under failure. Lemma \ref{lem:counting}(ii):
$(k-1) \cdot 9 \le 27$ gives $k \le 4$; $k \le 3$ is excluded by $g(3) = 7 < 8$ and
$g(2) = 5 < 8$; hence $k = 4$. This is the case $n$ even, $k = n/2$ of Lemma
\ref{lem:counting}(iii): $d = 1$ and $m_2 = m_3 = m_4 = 9$.
\end{proof}

\begin{theorem}\label{thm:n8}
Every set of $8$ points in the plane has at least two $8$-good distances.
\end{theorem}

We give two logically independent proofs. Route B (Section \ref{sec:n8B}) is
self-contained modulo the classification $E_7(3) = \{R_7, R_6^+\}$ of
\cite[Theorem~1]{ErFi96} and a finite exact computation reproduced in the ancillary file
\texttt{n8\_descent.py}; we present it as the main proof. Route A (Section
\ref{sec:n8A}) closes the case against Shinohara's classification of $8$-point
$4$-distance sets \cite{Shinohara08} and serves as an independent verification.

\subsection{Route B: descent to the heptagon (main proof)}\label{sec:n8B}

\begin{proposition}\label{prop:n8descent}
If $P$, $|P| = 8$, fails the conjecture, then $P = R_7 \cup \{v\}$ for a regular
heptagon $R_7$, where $v$ has exactly two vertices of $R_7$ at each of the three
distances of the heptagon (covering $6$ of its vertices), and the distance from $v$ to
the seventh vertex is a new value $\delta_1$ strictly larger than the heptagon's
diameter, realized only by that pair.
\end{proposition}

\begin{proof}
By Proposition \ref{prop:n8red}, $d = 1$, so (H2) holds automatically and Lemma
\ref{lem:descent} applies: deleting either endpoint $v$ of the unique diameter edge
leaves $K$, a set of $7$ points with exactly $3$ distances (each class survives with
$m_i' = 9 - a_i \ge 3$, where $a_i := \deg_i(v)$ and $\sum_{i \ge 2} a_i = 6$). There is
no size contradiction ($7 = g(3)$), but $K$ is a \emph{maximum} $7$-point $3$-distance
set, hence $K \in \{R_7, R_6^+\}$ by \cite[Theorem~1]{ErFi96}. Lemma
\ref{lem:descent}(d) gives $a_2 \ge 2$ and $m_i(K) = 9 - a_i \in [3, 9]$ with
$m_2(K) \le 7$.
\begin{itemize}
\item Base $R_6^+$, vector $(3, 6, 12)$: would require $9 - a_4 = 12$, i.e.\ $a_4 = -3 <
0$. Impossible by pure counting.
\item Base $R_7$, vector $(7,7,7)$: forces $a = (2,2,2)$ (consistent with $a_2 \ge 2$).
\end{itemize}
This is the statement.
\end{proof}

It remains to show that no such point $v$ exists. Normalize $R_7$ to circumradius $1$
with center $O$ (admissible: all constraints and the multiplicity profile are invariant
under similarity) and write $V_0, \dots, V_6$ for the vertices and
$D_s = 4\sin^2(s\pi/7)$, $s = 1, 2, 3$, for the three squared distances of $R_7$
($D_1 < D_2 < D_3$; $\delta_2^2 = D_3$).

\begin{lemma}[bisector rigidity]\label{lem:bisector}
Suppose $v \in \R^2$ has exactly two vertices of $R_7$ at each of the three squared
distances $D_1, D_2, D_3$. Then either
$v = O$, or $v$ lies on an axis of symmetry of $R_7$; in the latter case, if $V_w$ is
the vertex on that axis, the three equidistant pairs are the mirror pairs
$\{V_{w+1}, V_{w-1}\}$, $\{V_{w+2}, V_{w-2}\}$, $\{V_{w+3}, V_{w-3}\}$ (indices mod
$7$), and the leftover vertex is $V_w$.
\end{lemma}

\begin{proof}
If $|vP| = |vQ|$ with $P, Q$ on the unit circle centered at $O$, then $v$ lies on the
perpendicular bisector of $PQ$, which passes through $O$. The three equidistant pairs
give three such bisectors through $O$ all containing $v$; hence either $v = O$ or all
three bisectors are one and the same line $\ell \ni O$. The bisector of
$\{V_i, V_j\}$ is the line through $O$ with direction angle $\pi(i+j)/7 \pmod{\pi}$
(well defined: in a $7$-gon no chord is a diameter of the circumcircle), so two
bisectors coincide iff the index sums $i+j$ are congruent mod $7$. Among the $105$
partitions of $\Z_7$ into three pairs and a singleton, exactly $7$ have all three pair
sums congruent mod $7$ (mechanically verified, and by hand: if the singleton is $w$,
the total sum $21 \equiv 0$ forces the common pair sum to be $\equiv 2w \pmod 7$, and a
pair $\{w+a, w+b\}$ with $a + b \equiv 0 \pmod 7$, $a, b \in \pm\{1,2,3\}$, forces
$b = -a$): they are precisely the mirror partitions around the axis through $V_w$ and
$O$, one for each $w \in \Z_7$.
\end{proof}

\begin{proposition}[exact search]\label{prop:n8search}
No point $v \in \R^2 \setminus R_7$ has exactly two vertices of $R_7$ at each of the
three distances of the heptagon. Consequently the point $v$ of Proposition
\ref{prop:n8descent} does not exist.
\end{proposition}

\begin{proof}
The case $v = O$ of Lemma \ref{lem:bisector} is impossible: all seven distances from
$O$ equal the circumradius. (In passing, $R_7^+ = R_7 \cup \{O\}$ has profile
$(7,7,7,7)$, not $(1,9,9,9)$.) So $v = t\,u$ with $u$ the unit vector toward $V_w$,
$t \in \R$, and the pairs are the mirror pairs. The entire computation lives in the
totally real cubic field $K = \Q(y)$, $y = 2\cos(\pi/7)$, with minimal polynomial
$x^3 - x^2 - 2x + 1$: writing $z_m = 2\cos(m\pi/7) \in K$ (Chebyshev recurrence
$z_{m+1} = y z_m - z_{m-1}$), we have $D_s = 2 - z_{2s}$, explicitly
$D_1 = 4 - y^2$, $D_2 = y^2 - y + 1$, $D_3 = 2 + y$, with $0 < D_1 < D_2 < D_3$
certified exactly (signs in $K$ are decided by rational interval isolation of $y$, with
no floating point in any decision). The mirror pair $\{V_{w \pm r}\}$ lies at angular
distance $2\pi r/7$ from $V_w$, so assigning the squared distance $D_{\sigma(r)}$ to
that pair, for a bijection $\sigma$ of $\{1,2,3\}$, reads
\[
t^2 - 2\cos(2\pi r/7)\, t + 1 = D_{\sigma(r)}, \qquad r = 1, 2, 3 .
\]
The three cosines are distinct, so subtracting two of the equations determines
$t \in K$ uniquely; the remaining equation, and finally the checks on
$x := |v - V_w|^2 = (t-1)^2$ (namely $x \ne 0$, i.e.\ $v \notin R_7$, and
$x > D_3$ for a new diameter), are exact conditions in $K$. By the rotational symmetry
the $7$ axes behave identically, leaving $6$ assignments $\sigma$ to check. The result
of the exact enumeration (ancillary file \texttt{n8\_descent.py}): five of the six
assignments fail one of the three equalities with a nonzero residue in $K$ (nonzero
because its coordinates in the basis $\{1, y, y^2\}$ are not all zero --- an exact
criterion); the single assignment that satisfies all three equalities is the natural
one (pair $\{V_{w\pm r}\}$ at distance $D_r$), which is the expected positive control:
it forces $x = 0$, i.e.\ $v = V_w \in R_7$, so there is no eighth point. Survivors:
none.
\end{proof}

\begin{proof}[Proof of Theorem \ref{thm:n8} (Route B)]
If $P$ failed, Propositions \ref{prop:n8red} and \ref{prop:n8descent} would produce a
point $v \notin R_7$ with exactly two heptagon vertices at each heptagon distance,
contradicting Proposition \ref{prop:n8search}.
\end{proof}

\begin{remark}[redundant verifications]\label{rem:n8redundant}
The enumeration of Proposition \ref{prop:n8search} is cross-checked in three
independent ways (all in \texttt{n8\_descent.py}). (1) A second, structure-free
enumeration in the spirit of Corollary \ref{cor:extension}: every valid $v$ lies on the
intersection of two circles centered at distinct vertices with radii among the
$\sqrt{D_s}$; all $\binom{7}{2} \times 3^2 = 189$ circle pairs are enumerated
exactly over $K(\sqrt{\Delta})$ ($21$ pairs have no real intersection; of the $336$
real candidates, $231$ fail the required distance multiset and $105$ satisfy the
equalities but are degenerate). The $105$ degenerate candidates are \emph{exactly} the
seven vertices $V_w$, each recovered $\binom{6}{2} = 15$ times --- an exact internal
consistency check, and a proof that the machinery does find solutions when they exist.
(2) A Hopf--Pannwitz consistency check: a candidate passing the equalities with
$0 < x < D_3$ would yield a planar $8$-point set whose diameter has multiplicity
$7 + 2 = 9 > 8$, contradicting \cite{HoPa34}; coherently, no case fell in that branch.
(3) With \texttt{sympy} available, the full enumeration is re-run independently over
\texttt{QQ.algebraic\_field}$(2\cos(\pi/7))$, with the minimal polynomial re-derived
independently; the verdict ($35$ failures, $7$ degenerate cases, $0$ survivors) is
identical.
\end{remark}

\subsection{Route A: the classification of $8$-point $4$-distance sets (independent
verification)}\label{sec:n8A}

Shinohara \cite[Theorem~1.2(a), p.~3049]{Shinohara08} proves (as an intermediate result
toward the uniqueness of the maximum $12$-point $5$-distance set) that every $8$-point
four-distance set in $\R^2$ is similar to $R_8$, $R_7^+$, the configuration of his
Fig.~1(e), or an $8$-point subset of a $9$-point four-distance set. The relevance of
this classification to the profile $(1,9,9,9)$ was first observed in the forum comment
\cite{Chojecki26}. We have verified the statement against the full text and computed
the exact multiplicity vector of every class.

The $9$-point parents are the four configurations $R_9$, $B_9$, $A_9$, $T_9$ of
\cite[Theorem~1]{ErFi96} (see Section \ref{sec:n9} for their description and profiles).
Every $8$-point subset of a $9$-point $4$-distance set automatically has exactly $4$
distances (with $3$ distances it would be an $8$-point $3$-distance set, contradicting
$g(3) = 7$), so all $36$ vertex deletions are candidates; up to similarity they fall
into $12$ classes ($R_9$: one class; $B_9$: six; $A_9$: three; $T_9$: three; minus one
cross-parent coincidence: $B_9$ minus a far off-axis vertex is similar to $A_9$ minus a
spike, so that $8$-point set extends to a $9$-point $4$-distance set in two
inequivalent ways). The configuration of Fig.~1(e) is the unit square with the apex of
an equilateral triangle erected outward on each side (within the $D_4$-symmetric
one-parameter family, the four-distance condition forces the equilateral apex height
--- checked exactly; the paper gives the configuration only as a figure). Table
\ref{tab:e84} lists the resulting $15$ similarity classes with their exact
multiplicity vectors (computed over $\Q(\sqrt3)$ for the lattice-based configurations
and by orbit counting with exact chord-separation certificates for the regular
polygons; each vector sums to $28$; Hopf--Pannwitz $m_1 \le 8$ holds throughout).

\begin{table}[ht]
\centering
\caption{The $15$ similarity classes of $8$-point $4$-distance sets
(\cite[Theorem~1.2(a)]{Shinohara08} plus \cite[Theorem~1]{ErFi96}), with exact
multiplicity vectors ordered by decreasing distance. Parents: $B_9 = (4,7,9,16)$,
$A_9 = (6,6,9,15)$, $T_9 = (6,6,18,6)$, $R_9 = (9,9,9,9)$.}
\label{tab:e84}
\begin{tabular}{@{}rll@{}}
\toprule
\# & configuration & vector \\
\midrule
1 & $R_8$ (regular octagon; sporadic) & $(4,8,8,8)$ \\
2 & $R_7^+$ (regular heptagon with center; sporadic) & $(7,7,7,7)$ \\
3 & square with four equilateral apexes (Fig.~1(e); sporadic) & $(2,12,2,12)$ \\
4 & $R_9$ minus a vertex & $(7,7,7,7)$ \\
5 & $B_9$ minus the center & $(4,7,7,10)$ \\
6 & $B_9$ minus the near axial vertex & $(4,6,7,11)$ \\
7 & $B_9$ minus a vertex adjacent to one spike & $(4,5,7,12)$ \\
8 & $B_9$ minus a far off-axis vertex ($\cong A_9$ minus a spike) & $(3,5,7,13)$ \\
9 & $B_9$ minus the opposite axial vertex & $(2,6,7,13)$ \\
10 & $B_9$ minus a spike ($= R_6^+$ plus one spike) & $(2,5,7,14)$ \\
11 & $A_9$ minus a corner & $(5,4,7,12)$ \\
12 & $A_9$ minus an interior point & $(6,5,7,10)$ \\
13 & $T_9$ minus an inner vertex & $(6,5,14,3)$ \\
14 & $T_9$ minus an intermediate vertex & $(5,4,14,5)$ \\
15 & $T_9$ minus an outer vertex & $(3,5,14,6)$ \\
\bottomrule
\end{tabular}
\end{table}

Classes 2 and 4 have equal vectors but are not similar ($R_7^+$ has a point with seven
equal distances, $R_9$ minus a vertex has none; their chord fields are distinct cubic
fields). Classes 7 and 11 (and likewise 13 and 15) have equal multiplicity
\emph{multisets} but assign the diameter differently, so they are not similar either.

\begin{proof}[Proof of Theorem \ref{thm:n8} (Route A)]
If $P$ failed, it would be an $8$-point $4$-distance set with vector $(1,9,9,9)$
(Proposition \ref{prop:n8red}), hence similar to one of the $15$ classes of Table
\ref{tab:e84}. The multiplicity vector is a similarity invariant, and none of the $15$
vectors is $(1,9,9,9)$ --- not even as a multiset (no vector contains a $1$, and none
contains a $9$). Contradiction.
\end{proof}

\begin{remark}[comparison with the posted GPT-5.2 sketch]\label{rem:gpt}
Checked point by point against \cite{Shinohara08}, the sketch posted in
\cite{Chojecki26} matches the paper $1{:}1$ on every claim that can be checked: the
classification exists (Theorem 1.2(a), uncited in the sketch); the three sporadic
configurations are $R_8$, $R_7^+$ and ``the square with four equilateral bumps''
(which is exactly Fig.~1(e)); the profiles it stated spontaneously --- $R_8$ with
multiset $\{8,8,8,4\}$, $R_7^+$ with $(7,7,7,7)$, two classes of multiplicity $2$ for
the square with bumps --- are all correct. What the post was not is a verifiable
result: it cited no source, asserted ``multiplicity tables'' without exhibiting any,
and did not cover the deletions of the non-lattice parent $T_9$ --- whose profiles
$(6,5,14,3)$, $(5,4,14,5)$, $(3,5,14,6)$ had, to our knowledge, never been published
and appear for the first time in Table \ref{tab:e84}. The gap between ``substantially
correct'' and ``verified against the primary source in exact arithmetic'' is exactly
what this section closes.
\end{remark}

\section{The case $n = 9$}\label{sec:n9}

\begin{proposition}[reduction]\label{prop:n9red}
If $P$, $|P| = 9$, fails, then $P$ has exactly $4$ distinct distances,
$d \in \{2, \dots, 6\}$, and its diameter graph is not a star.
\end{proposition}

\begin{proof}
$\binom{9}{2} = 36$, $m_i \ge 10$. Lemma \ref{lem:counting}(ii): $(k-1) \cdot 10 \le 35$
gives $k \le 4$; $k \in \{2,3\}$ are excluded by $g(2) = 5 < 9$, $g(3) = 7 < 9$; hence
$k = 4$, and $d \le 36 - 30 = 6$. Corollary \ref{cor:size} with $g(3) = 7$: a star
diameter graph (in particular $d = 1$) would yield an $8$-point $3$-distance set,
impossible; hence $d \ge 2$ and no star.
\end{proof}

The possible vectors are $(d, 10+e_2, 10+e_3, 10+e_4)$ with $\sum e_i = 6 - d$,
$d \in \{2, \dots, 6\}$: $35$ vectors in total. None of them survives:

\begin{theorem}\label{thm:n9}
Every set of $9$ points in the plane has at least two $9$-good distances.
\end{theorem}

\begin{proof}
Suppose $P$, $|P| = 9$, fails. By Proposition \ref{prop:n9red}, $P$ is a $9$-point
$4$-distance set with $9 = g(4)$, hence maximum, and by the classification of
\cite[Theorem~1 and Fig.~1]{ErFi96} (verified against the full text) $P$ is similar to
one of exactly \emph{four} configurations:
\begin{itemize}
\item $R_9$, the regular nonagon;
\item $B_9 \subset \Lat$, the regular hexagon with center plus two adjacent spikes of
the hexagram;
\item $A_9 \subset \Lat$, the lattice piece with rows $2/3/3/1$;
\item $T_9$, three concentric equilateral triangles \emph{not} contained in $\Lat$
(inner and intermediate triangles pointing one way, outer triangle the other;
normalizing the inner circumradius to $1$, the circumradii are $1$, $1+\sqrt3$,
$2+\sqrt3$, and the four squared distances are $3$, $6+3\sqrt3$, $12+6\sqrt3$,
$21+12\sqrt3$).
\end{itemize}
Their exact multiplicity vectors, ordered by decreasing distance (ancillary file
\texttt{n9\_profiles.py}; Loeschian norms for $A_9, B_9$, the field $\Q(\sqrt3)$ for
$T_9$, orbit counting with chord monotonicity for $R_9$; each sums to $36$), are
\[
R_9 : (9,9,9,9), \qquad B_9 : (4,7,9,16), \qquad A_9 : (6,6,9,15), \qquad
T_9 : (6,6,18,6).
\]
The minimum non-diameter multiplicity is $9$, $7$, $6$, $6$ respectively --- in every
case $< 10$, so none matches the failure profile. Contradiction.
\end{proof}

\begin{remark}[independent support]\label{rem:n9support}
An exhaustive exact search over the triangular lattice returns exactly $\{A_9, B_9\}$
as the $9$-point $4$-distance subsets of $\Lat$ up to similarity, matching the two
lattice configurations of \cite[Fig.~1]{ErFi96}; the existence of $T_9$ and the fact
that it realizes exactly four distances are certified in $\Q(\sqrt3)$. What is not
independently re-derived here is the completeness of the non-lattice part of the
classification (that $R_9$ and $T_9$ are the only non-lattice examples); that rests on
the case analysis of \cite{ErFi96}, which the subsequent literature
(\cite{Shinohara08,Wei12,LanWei13}) treats as established. As a sanity check, the
multiplicity $18$ in $T_9$ attains Harborth's maximum for $9$ points and does not
exceed it. Note also that $B_9$ and $A_9$ show the failure window $d \in \{2,\dots,6\}$
of Proposition \ref{prop:n9red} is inhabited by real configurations ($d = 4$ and $6$);
they fail the profile only through their small classes.
\end{remark}

\section{The case $n = 10$}\label{sec:n10}

\begin{proposition}[reduction]\label{prop:n10red}
If $P$, $|P| = 10$, fails, then $P$ has exactly $5$ distinct distances with multiplicity
vector exactly $(1, 11, 11, 11, 11)$.
\end{proposition}

\begin{proof}
$\binom{10}{2} = 45$, $m_i \ge 11$. Lemma \ref{lem:counting}(ii): $(k-1) \cdot 11 \le 44$
gives $k \le 5$; $k \le 4$ is excluded by $g(4) = 9 < 10$ (and $g(3), g(2)$); hence
$k = 5$, the even rigid case of Lemma \ref{lem:counting}(iii): $d = 1$,
$m_i = 11$ for $i \ge 2$.
\end{proof}

\begin{theorem}\label{thm:n10}
Every set of $10$ points in the plane has at least two $10$-good distances.
\end{theorem}

We prove this \emph{unconditionally}, by descent to the completely classified $9$-point
$4$-distance sets of \cite[Theorem~1]{ErFi96} followed by a finite exact extension search.
The proof uses only $E_9(4)$ (a fully-proved classification, Section \ref{sec:n9}) and
Hopf--Pannwitz; it does \emph{not} depend on Wei's classification $E_{10}(5)$ \cite{Wei12}.
An independent corroboration against \cite[Theorem~11]{Wei12} --- which also yields an
erratum to \cite{Wei12} --- is recorded afterwards (Remark \ref{rem:n10wei}).

\emph{Descent to $E_9(4)$: a counting argument leaving only $R_9$.} Suppose $P$,
$|P| = 10$, fails; by Proposition \ref{prop:n10red} it is a $10$-point $5$-distance set
with the rigid vector $(1,11,11,11,11)$. Since $d = 1$, hypothesis (H2) holds
automatically and Lemma \ref{lem:descent} descends: deleting either endpoint $v$ of the
unique diameter edge leaves a $9$-point $4$-distance set $K = P \setminus \{v\}$ with
$m_i(K) = 11 - a_i$, where $a_i = \deg_i(v) \ge 0$ and $\sum_i a_i = 8$; thus
$m_i(K) \in \{3, \dots, 11\}$ for all four classes, and Hopf--Pannwitz on the $9$-point
set $K$ gives $m_1(K) \le 9$. Now $9 = g(4)$, so $K$ is a \emph{maximum} $9$-point
$4$-distance set, hence by \cite[Theorem~1]{ErFi96} (the classification of Section
\ref{sec:n9}, verified against the full text) $K$ is similar to $R_9$, $B_9$, $A_9$ or
$T_9$, with the exact multiplicity vectors
\[
R_9 : (9,9,9,9), \quad B_9 : (4,7,9,16), \quad A_9 : (6,6,9,15), \quad T_9 : (6,6,18,6)
\]
of Section \ref{sec:n9}. Adjoining one point can only \emph{add} edges to each class, so
$a_i \ge 0$ forces $m_i(K) \le 11$ for every class; but $B_9$, $A_9$ and $T_9$ each carry
a class of multiplicity $16$, $15$ and $18$ respectively, all $> 11$, so none can be the
base of a descent from $(1,11^4)$ --- they are excluded by pure counting, with no search.
Only $R_9$ survives: $(9,9,9,9)$ has every class $\le 11$, and the degrees of $v$ are then
forced with no freedom, $a_i = 11 - 9 = 2$ for each of the four chord classes
($\sum_i a_i = 8$, meeting $a_2 \ge 2$ of Lemma \ref{lem:descent}(d) with equality). Hence
$K \cong R_9$ and $v$ has exactly two vertices of $R_9$ at each of its four chord lengths
($8$ edges), plus one edge to its diameter partner $u \in K$ at a new value
$\delta_1 > \diam R_9$ (so $m_1(P) = 1$). This is exactly Corollary \ref{cor:extension}
with $K = R_9$, the degree-$4$ analogue of the closed cases $n = 8$ (Route B, Proposition
\ref{prop:n8search}: $K = R_7$, $a = (2,2,2)$) and $n = 12$, $k = 6$ (Proposition
\ref{prop:n12search}: $K = R_{11}$, $a = (2,2,2,2,2)$).

\begin{proposition}[exact extension search over $R_9$]\label{prop:n10search}
No point $v \in \R^2 \setminus R_9$ has exactly two vertices of the regular nonagon $R_9$
at each of its four chord lengths together with a single farther vertex $u \in R_9$ at a
new value $\delta_1 > \diam R_9$. Consequently no $10$-point set fails the conjecture.
\end{proposition}

\begin{proof}
This is the finite exact search of Corollary \ref{cor:extension} for $K = R_9$ with the
forced degree pattern $a = (2,2,2,2)$. It is carried out with no floating point, in the
degree-$3$ field $\Q(\zeta_9)^+ = \Q(\beta)$, $\beta = 2\cos(2\pi/9)$ (minimal polynomial
$t^3 - 3t + 1$; the four squared chord lengths of $R_9$ are
$4\sin^2(s\pi/9) = 2 - 2\cos(2\pi s/9)$, $s = 1, \dots, 4$, written exactly over $\beta$),
extended by the quadratic tower $\Q(\beta)(\sqrt{\Delta})$ that carries the
circle-intersection candidate. Up to the dihedral symmetry of $R_9$ the diameter partner
$u \in R_9$ ranges over $9$ choices, all enumerated; for each, the eight remaining vertices
are partitioned into four pairs, one at each chord length, and the resulting colouring pins
$v$ as an intersection of two circles centred at vertices of $R_9$ with radii among its
chord lengths. Each candidate is then filtered by exact necessary conditions: the closed
four-point Cayley--Menger determinants vanish (planarity), the characteristic polynomial of
the Gram matrix certifies it positive semidefinite of rank $\le 2$, the squared far
distance $z := |vu|^2$ solves its quadratic in the tower and strictly exceeds every squared
chord length of $R_9$ (a genuinely new diameter), and the final multiplicity profile equals
$(1,11,11,11,11)$. The enumeration (bin \texttt{e12\_descent
-{}-n10} of \texttt{rust/fewdist}, the same descent engine as Proposition
\ref{prop:n12search}) runs over all $9$ diameter partners, triggers $855$
closed-Cayley--Menger prunings and $9$ surviving $z$-candidates, and returns \emph{zero}
survivors. The engine is the one validated by the synthetic positive control of the
\texttt{-{}-self-check} analogue and by the planted extension recovered in the
\texttt{-{}-r11} run, so the empty output is a true non-existence rather than a vacuous one.
\end{proof}

\begin{proof}[Proof of Theorem \ref{thm:n10}]
If $P$ failed, the descent above would produce a point $v \notin R_9$ with exactly two
nonagon vertices at each of the four chord lengths of $R_9$ and one farther vertex at a new
diameter, contradicting Proposition \ref{prop:n10search}. The proof invokes only
\cite[Theorem~1]{ErFi96} (the classification $E_9(4)$, whose proofs are complete) and
Hopf--Pannwitz \cite{HoPa34}; hence $n = 10$ is closed unconditionally.
\end{proof}

\begin{remark}[independent corroboration via $E_{10}(5)$, and an erratum to
{\cite{Wei12}}]\label{rem:n10wei}
The same conclusion follows independently from Wei's classification of the $10$-point
$5$-distance sets, and that route uncovers an erratum in \cite{Wei12}. Wei
\cite[Theorem~11]{Wei12} classifies the $10$-point $5$-distance sets: every one is similar
to $R_9^+$, $R_{10}$, $R_{11}{-}1$ (the regular $11$-gon minus a vertex), the double $R_5$
of his Figure 2q (a regular pentagon together with the five intersection points of its
diagonals), or one of the sixteen configurations in $\Lat$ of his Figures 2a--2p. Up to
similarity this list comprises $19$ classes (not $20$: see Remark \ref{rem:weierratum});
their exact multiplicity vectors are given in Table \ref{tab:e105}. None is
$(1,11,11,11,11)$: the only classes with diameter multiplicity $1$ are those of
Figures 2a, 2c ($=$ 2h) and 2f, whose vectors $(1,6,8,11,19)$, $(1,7,8,11,18)$ and
$(1,6,9,11,18)$ each have at least two non-diameter entries $\le 9 < 11$. (Moreover
every lattice class has an entry $\ge 13$; $R_9^+$ and $R_{11}{-}1$ have constant
vector $(9,9,9,9,9)$; and all multiplicities of $R_{10}$ and of the double $R_5$ are
divisible by $5$, so they could never contain a $1$ or an $11$.) This reproves Theorem
\ref{thm:n10} through a wholly different classification, and along the way records the
count $19$ rather than Wei's $20$ (Remark \ref{rem:weierratum}).
\end{remark}

\begin{table}[ht]
\centering
\caption{The $19$ similarity classes of $10$-point $5$-distance sets
(\cite[Theorem~11]{Wei12}, corrected per Remark \ref{rem:weierratum}), with exact
multiplicity vectors ordered by decreasing distance. All lattice classes have squared
distances $9 > 7 > 4 > 3 > 1$ (Loeschian norms).}
\label{tab:e105}
\begin{tabular}{@{}llll@{}}
\toprule
Wei figure & vector & Wei figure & vector \\
\midrule
2a & $(1,6,8,11,19)$ & 2j & $(3,12,8,9,13)$ \\
2b & $(2,7,8,11,17)$ & 2k & $(3,10,8,9,15)$ \\
2c $\cong$ 2h & $(1,7,8,11,18)$ & 2l & $(2,9,8,10,16)$ \\
2d & $(2,6,8,11,18)$ & 2m & $(2,9,8,11,15)$ \\
2e & $(2,7,8,10,18)$ & 2n & $(2,9,9,10,15)$ \\
2f & $(1,6,9,11,18)$ & 2o & $(3,10,8,10,14)$ \\
2g ($=T_{10}$, side-3 triangle) & $(3,6,9,9,18)$ & 2p & $(2,8,8,10,17)$ \\
2i & $(3,8,9,9,16)$ & & \\
\midrule
$R_9^+$ & $(9,9,9,9,9)$ & $R_{11}{-}1$ & $(9,9,9,9,9)$ \\
$R_{10}$ & $(5,10,10,10,10)$ & double $R_5$ (2q) & $(5,5,15,15,5)$ \\
\bottomrule
\end{tabular}
\end{table}

\begin{remark}[erratum to {\cite[Theorem~11]{Wei12}}]\label{rem:weierratum}
Wei's Figures 2c and 2h depict mirror images of one and the same chiral similarity
class (verified exactly: their canonical forms agree up to reflection and differ
without it). Since ``isomorphic'' in \cite{Wei12} means similar --- and Wei's own list
identifies mirror images, as seven of the eight chiral lattice classes appear exactly
once --- the sixteen lattice figures represent fifteen similarity classes, and
$E_{10}(5)$ has $19$ classes, not $20$. No convention rescues the count of sixteen:
with reflections the lattice classes number $15$; without them they would number $23$.
\end{remark}

\begin{remark}[independent re-derivation of the lattice part of $E_{10}(5)$]\label{rem:weiverif}
The lattice part of Wei's theorem was re-derived here independently and exhaustively:
a search over all similarity classes of $10$-point subsets of $\Lat$ with exactly $5$
distances (exact integer arithmetic on Loeschian norms; canonical forms via primitive
reduction and the order-$12$ point group; a lemma showing lattice similarity classes
coincide with congruence classes of primitive representatives) finds exactly the $15$
classes of Table \ref{tab:e105}, stable for primitive diameter$^2 \le M$ with
$M = 30, 60, 100$; a second, prune-free brute-force enumeration ($94{,}143{,}280$
subsets) returns the identical list, and the structural split of Wei's proof (classes
containing a similar copy of $R_6^+$ versus not: $7 + 8$) is reproduced. The figures
2a--2q were extracted from the published PDF at 600\,dpi and matched exactly
(\texttt{extract\_figs.py}, \texttt{match\_figs.py}); the non-lattice profiles are
exact over $\Q(\sqrt5)$ and by orbit counting (\texttt{nonlattice\_exact.py}).
This re-derivation, the profiles of Table \ref{tab:e105} and the erratum of Remark
\ref{rem:weierratum} all belong to the \emph{corroboration} of Remark \ref{rem:n10wei}:
the main proof of Theorem \ref{thm:n10} (descent to $E_9(4)$ and the exact extension
search over $R_9$, Proposition \ref{prop:n10search}) does not invoke \cite{Wei12} at all,
so \emph{nothing in the case $n = 10$ rests on Wei's classification of $E_{10}(5)$}. In
particular the one point that the corroboration route would leave on Wei alone --- the
completeness statement that no non-lattice class exists beyond the four listed and no
lattice class has primitive diameter$^2 > 100$ (all known ones have diameter$^2 = 9$) ---
is no longer load-bearing for the theorem. Wei's proof in turn draws on
\cite{Shinohara04,Shinohara08,ErFi96,LanWei13}, on Lemmas 7--8 of \cite{Wei11}, and on the
diameter-graph analysis later published as \cite{WeiLiCongGao14}.
\end{remark}

\section{The case $n = 11$}\label{sec:n11}

\begin{proposition}[reduction]\label{prop:n11red}
If $P$, $|P| = 11$, fails, then $P$ has exactly $5$ distinct distances,
$d \in \{2, \dots, 7\}$, and its diameter graph is not a star.
\end{proposition}

\begin{proof}
$\binom{11}{2} = 55$, $m_i \ge 12$. Lemma \ref{lem:counting}(ii): $(k-1) \cdot 12 \le 54$
gives $k \le 5$; $k \le 4$ excluded by $g(4) = 9 < 11$ (and below); hence $k = 5$,
$d \le 55 - 48 = 7$. Corollary \ref{cor:size} with $g(4) = 9$: a star (in particular
$d = 1$) would yield a $10$-point $4$-distance set, impossible; so $d \ge 2$, no star.
\end{proof}

This is the reduction half of part (iii) of Theorem \ref{thm:main}. The possible vectors
$(d, 12+e_2, \dots, 12+e_5)$, $\sum e_i = 7 - d$, number $126$; abbreviate the resulting
constraint on the multiplicity vector (ordered by decreasing distance) of a hypothetical
$11$-point failure as
\[
\textbf{(T1)}\qquad m_1 = d \in \{2, \dots, 7\}, \qquad m_i \ge 12 \ \ (2 \le i \le 5).
\]
Here $11 < g(5) = 12$, so $P$ is not a maximum $5$-distance set and the descent of Section
\ref{sec:descent} produces no extension problem (it only excludes the star). What closes
the case is instead the classification of \emph{all} $11$-point $5$-distance sets.

\begin{theorem}[{Wei \cite[Theorem~13, p.~514]{Wei11}}]\label{thm:e115}
Every $11$-point $5$-distance set in the plane is similar to one of exactly four
configurations: the regular $11$-gon $R_{11}$, or one of the three triangular-lattice
patches $W_1, W_2, W_3 \subset \Lat$ of \cite[Figure~4, p.~515]{Wei11}.
\end{theorem}

In the embedding $(a,b) \mapsto a(1,0) + b(\tfrac12, \tfrac{\sqrt3}{2})$ of $\Lat$ (squared
distances the Loeschian norms $a^2+ab+b^2$), representatives of the three lattice classes
are
\begin{align*}
W_1 &: (0,0),(0,1),(0,2),(1,-1),(1,0),(1,1),(1,2),(2,-1),(2,0),(2,1),(3,-1); \\
W_2 &: (0,0),(0,1),(0,2),(1,-1),(1,0),(1,1),(1,2),(2,-1),(2,0),(3,-1),(3,0); \\
W_3 &: (0,0),(0,1),(0,2),(1,-1),(1,0),(1,1),(1,2),(2,-1),(2,1),(3,-1),(3,0).
\end{align*}
All three have squared-distance set $\{1,3,4,7,9\}$ and diameter$^2 = 9$; $W_1$ is chiral,
$W_2$ and $W_3$ are achiral. The remaining class $R_{11}$ is not contained in $\Lat$
(Remark \ref{rem:weiverif11}); we record its profile separately, since it is the only
class in the $n = 12$ window (Section \ref{sec:n12}) and its distances are algebraic of
degree $5$.

\begin{lemma}[profile of $R_{11}$]\label{lem:r11}
The regular $11$-gon determines exactly five distances, each realized by exactly $11$
pairs: its multiplicity vector, ordered by decreasing distance, is $(11,11,11,11,11)$,
and its diameter is the step-$5$ chord.
\end{lemma}

\begin{proof}
Label the vertices by $\Z/11$, $V_j = (\cos 2\pi j/11, \sin 2\pi j/11)$; the dihedral
group $D_{11}$ (rotation $r: j \mapsto j+1$, reflection $j \mapsto -j$) acts by isometries.
For $i \ne j$ the \emph{step} $s(\{i,j\}) := \min\bigl((i-j) \bmod 11, (j-i) \bmod 11\bigr)
\in \{1, \dots, 5\}$ is well defined ($11$ is odd, so the two residues are distinct), and
$|V_iV_j| = 2\sin(s\pi/11)$ depends only on $s$. The orbit of $\{0, s\}$ under the rotation
subgroup $C_{11}$ has $11$ elements, since $r^t\{0,s\} = \{0,s\}$ would force
$2s \equiv 0 \pmod{11}$, impossible for $1 \le s \le 5$; so each of the five step classes
contains exactly $11$ pairs, and $5 \cdot 11 = 55 = \binom{11}{2}$ accounts for every
pair. The five values $2\sin(s\pi/11)$, $s = 1, \dots, 5$, are strictly increasing (as
$5\pi/11 < \pi/2$ and $\sin$ is increasing on $[0,\pi/2]$), so there are exactly five
distinct distances and the diameter is the step-$5$ chord. Algebraically, with
circumradius $1$ the five squared distances $2 - 2\cos(2\pi s/11)$ are the roots of
$u^5 - 11u^4 + 44u^3 - 77u^2 + 55u - 11$, obtained from the minimal polynomial
$y^5 + y^4 - 4y^3 - 3y^2 + 3y + 1$ of $2\cos(2\pi/11)$ by the substitution $y = 2 - u$;
this quintic is Eisenstein at $11$, hence irreducible and separable, confirming that the
five squared distances are distinct and generate the real subfield $\Q(\zeta_{11})^+$ of
degree $5$ over $\Q$.
\end{proof}

\begin{table}[ht]
\centering
\caption{The four similarity classes of $E_{11}(5)$ (\cite[Theorem~13]{Wei11}), with
exact multiplicity vectors ordered by decreasing distance (each summing to $55$; ancillary
file \texttt{n11\_profiles.py}). The three lattice classes have squared distances
$9 > 7 > 4 > 3 > 1$ (Loeschian norms) and diameter$^2 = 9$; $R_{11}$ has the five squared
chord lengths $2 - 2\cos(2\pi s/11)$, $s = 1, \dots, 5$ (Lemma \ref{lem:r11}).}
\label{tab:e115}
\begin{tabular}{@{}llll@{}}
\toprule
class & Wei Figure~4 & vector & chiral \\
\midrule
$R_{11}$ & --- (non-lattice) & $(11,11,11,11,11)$ & no \\
$W_1$ & right (zigzag with notch) & $(2,9,10,13,21)$ & yes \\
$W_2$ & centre (trapezoid) & $(3,10,10,12,20)$ & no \\
$W_3$ & left (hexagon with hole) & $(3,12,10,12,18)$ & no \\
\bottomrule
\end{tabular}
\end{table}

\begin{theorem}\label{thm:n11}
Conditionally on the classification $E_{11}(5)$ of Wei \cite[Theorem~13]{Wei11}
(Theorem \ref{thm:e115}), every set of $11$ points in the plane has at least two
$11$-good distances.
\end{theorem}

\begin{proof}
Suppose $P$, $|P| = 11$, fails. By Proposition \ref{prop:n11red} its multiplicity vector
satisfies (T1) and $P$ is an $11$-point $5$-distance set, hence by Theorem \ref{thm:e115}
similar to $R_{11}$, $W_1$, $W_2$ or $W_3$. The multiplicity vector is a similarity
invariant, and the four vectors (Table \ref{tab:e115}) are $(11,11,11,11,11)$,
$(2,9,10,13,21)$, $(3,10,10,12,20)$ and $(3,12,10,12,18)$; none satisfies (T1). Indeed
each lattice class has a non-diameter entry $\le 10 < 12$ ($W_1$ has $m_2 = 9$ and
$m_3 = 10$; $W_2$ has $m_2 = m_3 = 10$; $W_3$ has $m_3 = 10$), while $R_{11}$ has
$m_1 = 11 > 7$ and every entry equal to $11 < 12$. Contradiction.
\end{proof}

This proves part (iii) of Theorem \ref{thm:main}.

\begin{remark}[independent verification of the lattice part, and what rests on
\cite{Wei11}]\label{rem:weiverif11}
The lattice part of Theorem \ref{thm:e115} was re-derived here independently and
exhaustively up to primitive diameter$^2 \le 300$: a search over all similarity classes of $11$-point subsets of $\Lat$ with
exactly $5$ distances (exact integer arithmetic on Loeschian norms; canonical forms via
primitive Eisenstein reduction and the order-$12$ point group, as in Section
\ref{sec:n10}) returns exactly the three classes $W_1, W_2, W_3$ of Table \ref{tab:e115},
stable for primitive diameter$^2 \le M$ with $M = 30, 60, 100, 147, 300$; the same
enumerator reproduces the fifteen lattice classes of $E_{10}(5)$ as a sanity check, and
each of $W_1, W_2, W_3$ matches a figure of \cite[Figure~4]{Wei11} by exact canonical form
--- a clean bijection, with no erratum of the kind found for $E_{10}(5)$ (Remark
\ref{rem:weierratum}). The four multiplicity profiles are recomputed in exact arithmetic
(ancillary file \texttt{n11\_profiles.py}; Loeschian norms for $W_1, W_2, W_3$, the orbit
argument of Lemma \ref{lem:r11} for $R_{11}$), each summing to $55$; and $R_{11}$ is
similar to no lattice set, since in $\Lat$ every ratio of squared distances is rational
whereas $\delta_2^2/\delta_1^2 = 2 + 2\cos(2\pi/11)$ for $R_{11}$ is irrational of degree
$5$ over $\Q$. What still rests on \cite{Wei11} alone is the completeness of the
non-lattice part --- that $R_{11}$ is the only similarity class outside $\Lat$ --- and the
absence of lattice classes with primitive diameter$^2 > 300$ (all three known ones have
diameter$^2 = 9$); such a full plane-wide re-derivation of $E_{11}(5)$, in the exact style
of \cite{SzOs20}, would in principle remove the conditional and appears feasible, but we
have not carried it out --- our exhaustive re-derivation covers the triangular-lattice
part only.
\end{remark}

\begin{remark}[on the source \cite{Wei11}]\label{rem:weiomit}
Unlike the other classification inputs of this note, \cite{Wei11} states at the outset
that ``in the following some proofs are omitted because of the restriction of the length
of the paper'' \cite[p.~506]{Wei11}, and several steps of its case analysis (the end of
its Lemma~7, the non-$4P_2$ cases of its Lemma~8, and parts of its Theorem~13) are only
sketched. This is the same editorial style as \cite{Wei12}, in which the present note
found a genuine erratum (Remark \ref{rem:weierratum}); it is why $n = 11$ is stated
conditionally on \cite{Wei11} (Theorem \ref{thm:n11}) and accompanied by the independent
verification of Remark \ref{rem:weiverif11}. We found no error in \cite{Wei11}: the three
lattice classes it lists are exactly the three we re-derive.
\end{remark}

\section{The case $n = 12$}\label{sec:n12}

\begin{proposition}[reduction]\label{prop:n12red}
If $P$, $|P| = 12$, fails, then either
\begin{itemize}
\item[(a)] $k = 6$ and the multiplicity vector is exactly $(1, 13, 13, 13, 13, 13)$, or
\item[(b)] $k = 5$, $d \in \{2, \dots, 12\}$, and the diameter graph is not a star.
\end{itemize}
\end{proposition}

\begin{proof}
$\binom{12}{2} = 66$, $m_i \ge 13$. Lemma \ref{lem:counting}(ii): $(k-1) \cdot 13 \le 65$
gives $k \le 6$; $k \le 4$ excluded by $g(4) = 9 < 12$ (and below); so $k \in \{5, 6\}$.
If $k = 6$: even rigid case, $d = 1$, all $m_i = 13$ --- branch (a). If $k = 5$:
$d \le 66 - 52 = 14$, cut to $d \le 12$ by Hopf--Pannwitz; Corollary \ref{cor:size} with
$g(4) = 9$ excludes the star (an $11$-point $4$-distance set would result), so
$d \ge 2$ --- branch (b).
\end{proof}

\begin{proposition}[branch (b) is impossible]\label{prop:n12k5}
No $P$ with $|P| = 12$ fails the conjecture with $k = 5$.
\end{proposition}

\begin{proof}
In branch (b), $|P| = 12 = g(5)$ (proved in \cite[Section~4]{ErFi96}), so $P$ is a
maximum $5$-distance set. By the uniqueness theorem of Shinohara
\cite[Theorem~1.2(b), p.~3049]{Shinohara08} --- verbatim: ``The configuration given in
Fig.~1(d) is the only 12-point five-distance set in $\R^2$'' (proof in his \S4(II),
p.~3054; verified against the full text) --- $P$ is similar to the triangular-lattice
patch
\begin{align*}
T_{12} = {}& \{(0,2),(1,2)\} \cup \{(0,1),(1,1),(2,1)\} \\
& {} \cup \{(0,0),(1,0),(2,0),(3,0)\} \cup \{(1,-1),(2,-1),(3,-1)\} \;\subset\; \Lat
\end{align*}
(rows $2/3/4/3$). Its exact multiplicity vector, computed over the Loeschian norms
(squared distances $9 > 7 > 4 > 3 > 1$ with multiplicities $3, 12, 12, 15, 24$, summing
to $66$), is
\[
T_{12} : (3, 12, 12, 15, 24) \quad\text{ordered by decreasing distance.}
\]
Its second and third classes have multiplicity $12 < 13$, so $T_{12}$ does not match
the failure profile (indeed it has three $12$-good distances). Contradiction.
\end{proof}

\begin{remark}
$T_{12}$ shows the threshold ``$\le n$'' in the problem statement is tight at $n = 12$:
two of its distances have multiplicity exactly $n$. Historical note: Erd\H{o}s and
Fishburn proved $g(5) = 12$ but explicitly left the uniqueness of the $12$-set open
(\cite[p.~123]{ErFi96}: ``we presently have no guarantee that the 12-set of Fig.~1 is
the only realizer of $g(5)=12$''); Shinohara's Theorem 1.2(b) supplies exactly that
piece, and its proof is based on his classification of $8$-point $4$-distance sets ---
the same input as Route A of Section \ref{sec:n8A}.
\end{remark}

Together with Proposition \ref{prop:n12red}, Proposition \ref{prop:n12k5} proves part
(ii) of Theorem \ref{thm:main}: any failure at $n = 12$ lives in branch (a), with the
rigid vector $(1, 13^5)$.

\emph{Branch (a), $k = 6$: a finite exact extension search over $R_{11}$, with no
solution.} With $d = 1$, (H2) is automatic and Lemma \ref{lem:descent} descends: deleting either endpoint
$v$ of the unique diameter edge leaves an $11$-point $5$-distance set $K = P \setminus
\{v\}$ with $m_i(K) = 13 - a_i$, where $a_i = \deg_i(v) \ge 0$ and $\sum_i a_i = 10$; thus
$m_i(K) \in \{3, \dots, 13\}$ for all five classes, and Hopf--Pannwitz on the $11$-point
set $K$ gives $m_1(K) \le 11$. Abbreviate this window on the profile of $K$ (ordered by
decreasing distance) as
\[
\textbf{(T2)}\qquad m_i(K) \in \{3, \dots, 13\} \ \text{for all } i, \qquad m_1(K) \in
\{3, \dots, 11\}.
\]
By Theorem \ref{thm:e115}, $K$ is similar to $R_{11}$, $W_1$, $W_2$ or $W_3$; comparing
with the profiles of Table \ref{tab:e115}, each of the three lattice classes has a class
of multiplicity $21$, $20$ or $18$, all $> 13$, so \emph{only $R_{11}$ satisfies (T2)}.
Hence $K \cong R_{11}$, and the degrees of $v$ are then forced with no counting freedom:
$a_i = 13 - m_i(R_{11}) = 13 - 11 = 2$ for every class, i.e.\ $v$ has exactly two vertices
of $R_{11}$ at each of its five chord lengths ($10$ edges), plus one edge to its diameter
partner $u \in K$ at a new value $\delta_1 > \diam R_{11}$ (so $m_1(P) = 1$); the value
$a_2 = 2$ meets the bound $a_2 \ge 2$ of Lemma \ref{lem:descent}(d) with equality. This is
exactly the situation of Corollary \ref{cor:extension} with $K = R_{11}$, and the
degree-$5$ analogue of the closed case $n = 8$ (Route B, Proposition \ref{prop:n8search},
where $K = R_7$ and $a = (2,2,2)$): the candidates $v$ are the finitely many intersections
of pairs of circles centred at vertices of $R_{11}$ with radii among its five chord
lengths, and every remaining check is an exact identity or inequality in the degree-$5$
field $\Q(\zeta_{11})^+$ (Lemma \ref{lem:r11}). Note that $E_{12}(6)$ itself is \emph{not}
classified in the literature; it is the descent, together with the profile analysis of
Section \ref{sec:n11}, that reduces this branch to a single $11$-point configuration.

\begin{proposition}[exact extension search over $R_{11}$]\label{prop:n12search}
No point $v \in \R^2 \setminus R_{11}$ has exactly two vertices of $R_{11}$ at each of its
five chord lengths together with a single farther vertex $u$ at a new value
$\delta_1 > \diam R_{11}$. Consequently branch \emph{(a)} is empty: no $12$-point set
fails the conjecture with $k = 6$.
\end{proposition}

\begin{proof}
This is the finite exact search of Corollary \ref{cor:extension} for $K = R_{11}$ with the
forced degree pattern $a = (2,2,2,2,2)$. It is carried out with no floating point, in the
degree-$5$ field $\Q(\zeta_{11})^+ = \Q(\alpha)$, $\alpha = 2\cos(2\pi/11)$ (minimal
polynomial $y^5 + y^4 - 4y^3 - 3y^2 + 3y + 1$, Lemma \ref{lem:r11}), extended by the
quadratic tower $\Q(\alpha)(\sqrt{\Delta})$ that carries the circle-intersection candidate.
Up to the dihedral symmetry of $R_{11}$ the diameter partner $u \in R_{11}$ ranges over
$11$ choices, all enumerated; for each, the ten remaining vertices are partitioned into
five pairs, one at each chord length, and the resulting colouring pins $v$ as an
intersection of two circles centred at vertices of $R_{11}$ with radii among its chord
lengths. Each candidate is then filtered by exact necessary conditions: the closed
four-point Cayley--Menger determinants vanish (planarity), the characteristic polynomial of
the Gram matrix certifies it positive semidefinite of rank $\le 2$, the squared far
distance $z := |vu|^2$ solves its quadratic in the tower and satisfies $z$ strictly greater
than every squared chord length of $R_{11}$ (a genuinely new diameter), and the final
multiplicity profile equals $(1, 13^5)$. The enumeration (bin \texttt{e12\_descent -{}-r11}
of \texttt{rust/fewdist}) runs over all $11$ diameter partners, triggers $2310$
closed-Cayley--Menger prunings, and returns \emph{zero} survivors. As a
positive control, the same search recovers a planted synthetic extension, so the empty
output is a true non-existence rather than a vacuous one.
\end{proof}

Together with Proposition \ref{prop:n12k5}, this closes both branches and proves part (ii)
of Theorem \ref{thm:main}: conditionally on \cite{Wei11}, no $12$-point set fails the
conjecture. The classification input $E_{11}(5)$ is now in hand (Section \ref{sec:n11}), so
the last remaining dependence is on the completeness of \cite[Theorem~13]{Wei11}, which a
full plane-wide re-derivation of $E_{11}(5)$ (Remark \ref{rem:weiverif11}) would also
discharge if carried out.

\section{The case $n = 13$}\label{sec:n13}

\begin{theorem}\label{thm:n13}
Every set of $13$ points in the plane has at least two $13$-good distances.
\end{theorem}

\begin{proof}
Suppose $P$, $|P| = 13$, fails; $\binom{13}{2} = 78$, $m_i \ge 14$. Lemma
\ref{lem:counting}(ii): $(k-1) \cdot 14 \le 77$ gives $k - 1 \le 5.5$, so $k \le 6$; and
$k \le 5$ is excluded by $g(5) = 12 < 13$ \cite{ErFi96}. Hence $k = 6$ and $P$ is a
$13$-point $6$-distance set with $13 = g(6)$ \cite{Wei12}, i.e.\ maximum. By the
classification of Sz\"oll\H{o}si and \"Osterg\aa rd \cite[Theorem~16]{SzOs20} (in
agreement with \cite{Wei12}), there are exactly three such configurations up to
isometry (of the normalized sets; i.e.\ up to similarity): the regular $13$-gon
$R_{13}$, the regular $12$-gon with center $R_{12}^+$, and the regular hexagram with
center. Their exact multiplicity vectors, ordered by decreasing distance (ancillary
file \texttt{n13\_profiles.py}; each sums to $78$), are
\[
R_{13} : (13, 13, 13, 13, 13, 13), \qquad
R_{12}^+ : (6, 12, 12, 12, 24, 12), \qquad
\text{hexagram}^+ : (3, 6, 12, 15, 18, 24).
\]
In detail: $R_{13}$ has $13$ pairs in each chord class, the chords
$2\sin(j\pi/13)$ being strictly increasing in $j$. For $R_{12}^+$ with circumradius
$1$ the squared distances are $4 > 2+\sqrt3 > 3 > 2 > 1 > 2-\sqrt3$; the center-vertex
distance $1$ coincides exactly with the step-$2$ chord ($|V_0V_2|^2 = 1$), which is the
structural coincidence keeping $k = 6$ and giving that class $12 + 12 = 24$ pairs. The
hexagram with center lies in $\Lat$ (hexagon of norm $1$, spikes of norm $3$, center),
with squared distances $12 > 9 > 7 > 4 > 3 > 1$ and multiplicities
$3, 6, 12, 15, 18, 24$. The diameter multiplicities $13, 6, 3$ all satisfy
Hopf--Pannwitz. The minimum non-diameter multiplicities are $13$, $12$ and $6$
respectively --- in every case $< 14$, so none matches the failure profile (indeed the
three configurations have $6$, $5$ and $3$ thirteen-good distances). Contradiction.
\end{proof}

\begin{remark}
We cite the journal numbering of \cite{SzOs20}: their Theorem 16 is Theorem 4.6 of the
arXiv version (1804.06040); the two carry the same classification statement.
The classification is computer-assisted (isomorph-free exhaustive generation plus
Gr\"obner-basis computations); we did not re-derive it, but the three profiles above
are re-derived here in exact arithmetic (integer orbit counting for $R_{13}$;
$\Z[\sqrt3]$ for $R_{12}^+$; Loeschian norms for the hexagram), confirming the
$g(6) = 13$ landscape of \cite{Wei12}.
\end{remark}

\section{Concluding remarks}\label{sec:remarks}

\begin{itemize}
\item \textbf{Status of the range.} Every case $7 \le n \le 13$ is now closed. Two of them
are settled by descent to a classified maximum set followed by a finite exact extension
search: $n = 10$, closed \emph{unconditionally} by descent to $E_9(4)$ of \cite{ErFi96}
(only the regular nonagon $R_9$ survives the counting, and the extension search over $R_9$
with the forced pattern $a = (2,2,2,2)$ over $\Q(\zeta_9)^+$ returns no solution,
Proposition \ref{prop:n10search}, Section \ref{sec:n10}); and the $k = 6$ branch of
$n = 12$ (Theorem \ref{thm:main}(ii)), the extension search adjoining one point to the
regular $11$-gon $R_{11}$ with the forced pattern $a = (2,2,2,2,2)$ over $\Q(\zeta_{11})^+$,
likewise with no solution (Proposition \ref{prop:n12search}, Section \ref{sec:n12}). What
is \emph{not} self-contained is the classification $E_{11}(5)$ of \cite{Wei11}, on which
$n = 11$ (Theorem \ref{thm:main}(iii)) and $n = 12$ both rest \emph{conditionally}, and
whose printed proofs are partly omitted for length (Remark \ref{rem:weiomit}); its
triangular-lattice part we have re-derived exactly (Remark \ref{rem:weiverif11}). The case
$n = 10$ no longer depends on Wei at all (its old route via $E_{10}(5)$ \cite{Wei12} is now
kept only as an independent corroboration, Remark \ref{rem:n10wei}). A full independent
re-derivation of $E_{11}(5)$ in the plane, in the exact style of \cite{SzOs20}, would make
both $n = 11$ and $n = 12$ unconditional.
\item \textbf{$n \ge 14$.} For $n = 14$ the counting forces $k = 7$, and $7$-distance
sets are not classified (\cite{SzOs20} describe such a classification as feasible ``in
the near-term future''); this is the current hard frontier of the finite-case method.
\item \textbf{Asymptotics.} Nothing in this note touches the second question of Erd\H{o}s
and Pach (the number of $n$-good distances tending to infinity); by \cite{GuthKatz15} the
counting reduction is void for large $n$. As the problem page notes, the general question
``cannot be resolved with a finite computation''; the finite cases are partial results of
exactly the kind Erd\H{o}s rewarded for $n = 5, 6$ \cite{ErFi95}.
\item \textbf{Multiplicity constructions.} Bhowmick \cite{Bhowmick24} constructs $n$-point
sets with $\lfloor n/4 \rfloor$ distances each of multiplicity $> n$; the failure profile
of Table \ref{tab:profiles} would require about $n/2 - 1$ such rich classes, twice what is
known to be constructible --- consistent with, but far from proving, the conjecture.
\item \textbf{Formalization.} The statement of the problem and of the variants studied
here is formalized in Lean~4 in the \textsc{Formal Conjectures} project
\cite{FormalConjectures}; the proofs of the finite cases are not formalized, as they rest
on classification results outside the current Lean libraries.
\end{itemize}

\subsection*{Acknowledgements}
The author thanks R.\ Zeraoulia, whose working note \cite{Zeraoulia26} first posted the
counting reduction, the case $n = 7$, and the reduction of $n = 8$ to the profile
$(1,9,9,9)$; P.\ Chojecki, whose forum comment \cite{Chojecki26} pointed to the
classification of planar $8$-point $4$-distance sets as closing $n = 8$; and T.\ F.\ Bloom
for maintaining the Erd\H{o}s problems database \cite{Bloom132}.

The proofs, the exact-arithmetic computations (the accompanying \texttt{Rust} and
\texttt{Python} code), and the exposition of this note were developed with substantial
assistance from Anthropic's Claude language models. The author directed the work, verified
the results, and is solely responsible for any remaining errors.

\bibliographystyle{plain}
\bibliography{refs}

\end{document}
